\documentclass[11pt,a4paper]{article}
\usepackage[margin=27mm]{geometry}
\usepackage[T1]{fontenc}\usepackage[utf8]{inputenc}\usepackage{lmodern}
\usepackage{xcolor,hyperref}
\definecolor{terra}{HTML}{8B4A3A}\hypersetup{colorlinks=true,urlcolor=terra}
\pagestyle{empty}\setlength{\parindent}{0pt}
\newcommand{\motivo}[1]{\vspace{6pt}\textbf{\color{terra}#1}\par}
\begin{document}
\begin{flushright}
\textbf{Giulia Bianchi}\\
Via delle Rose 18\\50123 Firenze\\
\href{mailto:giulia.bianchi@example.org}{giulia.bianchi@example.org}
\end{flushright}
\vspace{8pt}\color{terra}\rule{\linewidth}{1pt}\color{black}

Firenze, 18 settembre 2026

Alla cortese attenzione della Commissione di selezione\\
Centro Studi Mediterranei\\Via Roma 40\\00100 Roma

\motivo{Oggetto: candidatura per il ruolo di ricercatrice}

Gentili membri della Commissione,

desidero sottoporre la mia candidatura per il ruolo di ricercatrice presso il
Centro Studi Mediterranei. Il mio lavoro unisce ricerca archivistica,
comunicazione chiara e strumenti digitali aperti per rendere le fonti più
facilmente consultabili e confrontabili.

Nel mio incarico attuale ho coordinato un piccolo progetto di catalogazione,
preparato materiali per studenti e collaborato con ricercatori di discipline
diverse. Mi interessa in particolare il modo in cui il Centro collega la
ricerca locale a un pubblico internazionale.

\motivo{Esperienze rilevanti}
\begin{itemize}
\item ricerca e gestione di collezioni digitali;
\item redazione di testi in italiano e inglese;
\item organizzazione di workshop e documentazione riproducibile.
\end{itemize}

Ringraziandovi per l'attenzione, resto a disposizione per un colloquio.

Cordiali saluti,\\[1.2em]\textbf{Giulia Bianchi}

\vfill\small\textcolor{terra}{Allegato: curriculum vitae.}
\end{document}
